\textbf{Критерий согласия Пирсона (критерий хи-квадрат)}

Пусть выборка $\displaystyle X_{1} ,\ \dotsc ,\ X_{n}$ с распределением $\displaystyle P( X_{1} =a_{i}) =p_{i} ,\ i=\overline{1,m} ,\ \sum _{i=1}^{m} p_{i} =1$ (полиномиальная схема). Положим $\displaystyle \nu _{j} =\sum _{i=1}^{n} I_{X_{i} =a_{j}} ,\ j=\overline{1,m}$ -- количество осуществлений исхода $\displaystyle a_{j}$. Тогда $\displaystyle \sum _{i=1}^{m} \nu _{i} =n$. Вектор $\displaystyle \overline{p} =\begin{pmatrix}
p_{1}\\
\dotsc \\
p_{m}
\end{pmatrix}$ считается неизвестным. Хотим проверить простую гипотезу $\displaystyle H_{0} :\ \overline{p} =\overline{p}_{0}$, где $\displaystyle \overline{p}_{0} =\begin{pmatrix}
p_{1}^{0}\\
\dotsc \\
p_{m}^{0}
\end{pmatrix}$ -- заданный вектор с положительными координатами, и $\displaystyle \sum _{i=1}^{m} p_{i}^{0} =1$.
\begin{definition}
    \textit{Статистикой хи-квадрат Пирсона} называют величину
    \begin{equation*}
        \hat{\chi }_{n}^{2} =\sum _{j=1}^{m}\dfrac{\left( \nu _{j} -np_{j}^{0}\right)^{2}}{np_{j}^{0}} .
    \end{equation*}
\end{definition}
\begin{theorem}
    (Пирсона) Если $\displaystyle H_{0}$ верна, то $\displaystyle \hat{\chi }_{n}^{2}\xrightarrow[n\rightarrow \infty ]{d} \chi _{m-1}^{2}$.
\end{theorem}

\begin{definition}
    Пусть $\displaystyle u_{1-\varepsilon }$ -- $\displaystyle ( 1-\varepsilon )$-квантиль $\displaystyle \chi _{m-1}^{2}$. Тогда \textit{критерием хи-кварат называется} $\displaystyle \left\{\hat{\chi }_{n}^{2}  >u_{1-\varepsilon }\right\}$.
\end{definition}
\begin{note}
    На практике критерий применяется при $\displaystyle np_{j}^{0} \geqslant 5,\ 1\leqslant j\leqslant m$.
\end{note}
\begin{note}
    Критерий хи-квадрат является асимптотическим, т.е. вероятность ошибки только стремится к $\displaystyle \varepsilon $ при $\displaystyle n\rightarrow \infty $.
\end{note}
\begin{proposition}
    Критерий хи-квадрат состоятелен.
\end{proposition}
\begin{proof}
    Пусть на самом деле $\displaystyle \overline{p} \neq \overline{p}_{0}$. Тогда $\displaystyle \hat{\chi }_{n}^{2} =n\sum _{i=1}^{m}\left(\dfrac{\nu _{i}}{n} -p_{i}^{0}\right)^{2}\dfrac{1}{p_{i}^{0}}$. В силу УЗБЧ для любого фиксированного $\displaystyle i\in \{1,\ \dotsc ,\ m\}$ выполнено $\displaystyle \dfrac{\nu _{i}}{n} =\dfrac{1}{n}\sum _{j=1}^{n} I_{X_{j} =a_{i}}\xrightarrow{a.s.} p_{i}$. По теореме о наследовании сходимости $\displaystyle \sum _{i=1}^{m}\left(\dfrac{\nu _{i}}{n} -p_{i}^{0}\right)^{2} \cdotp \dfrac{1}{p_{i}^{0}}\xrightarrow{a.s.}\sum _{i=1}^{m}\dfrac{\left( p_{i} -p_{i}^{0}\right)^{2}}{p_{i}^{0}}  >0$. Значит, $\displaystyle \hat{\chi }_{n}^{2}\xrightarrow[n\rightarrow \infty ]{} \infty $, а, следовательно, и $\displaystyle P\left(\hat{\chi }_{n}^{2}  >u_{1-\varepsilon }\right)\xrightarrow[n\rightarrow \infty ]{} 1$. Таким образом, вероятность ошибки второго рода стремится к нулю.
\end{proof}


\textbf{Критерий согласия Колмогорова (Колмогорова-Смирнова)}


Пусть $\displaystyle X=( X_{1} ,\ \dotsc ,\ X_{n})$ -- выборка из неизвестного распределения с непрерывной функцией распределения $\displaystyle F$. Тогда $\displaystyle D_{n} =\sup _{x\in \mathbb{R}}\left| F_{n}^{*}( x) -F( x)\right| $.
\begin{theorem}
    (Колмогорова, б/д) Пусть $\displaystyle F$ -- непрерывна. Тогда
    
    1) Распределение $\displaystyle \sqrt{n} D_{n}$ не зависит от вида $\displaystyle F$
    
    2) $\displaystyle \sqrt{n} D_{n}\xrightarrow{d} \xi $, где $\displaystyle \xi $ имеет распределение Колмогорова, т.е. $\displaystyle P( \xi \leqslant z) =\sum _{j\in \mathbb{Z}}( -1)^{j} e^{-2j^{2} z^{2}} I_{z >0}$.
\end{theorem}
Рассмотрим гипотезу $\displaystyle H_{0} :\ F=F_{0}$. - непрерывна
\begin{definition}
    \textit{Критерием Колмогорова} называется
    \begin{equation*}
        S=\left\{\sqrt{n}\sup _{x\in \mathbb{R}}\left| F_{n}^{*}( x) -F_{0}( x)\right|  >K_{1-\alpha }\right\} ,
    \end{equation*}
    где $\displaystyle K_{1-\alpha }$ -- $\displaystyle ( 1-\alpha )$-квантиль распределения Колмогорова.
\end{definition}
\begin{note}
    На практике критерий применяется при $\displaystyle n\geqslant 20$.
\end{note}


\textbf{Критерий омега-квадрат (или критерий фон Мизеса-Смирнова)}

Пусть гипотеза $\displaystyle H_{0} :\ F=F_{0}$, где $\displaystyle F_{0}$ -- непрерывная функция распределения.
\begin{definition}
    \textit{Статистикой омега-квадрат} называется
    \begin{equation*}
        \omega ^{2} :=n\int \left| F_{n}^{*}( x) -F_{0}( x)\right| ^{2} dF_{0}( x) .
    \end{equation*}
\end{definition}
\begin{proposition} (б/д)
    Если гипотеза $\displaystyle H_{0}$ верна, то статистика $\displaystyle \omega ^{2}$ слабо сходится к известному распределению.
\end{proposition}
\begin{definition}
    \textit{Критерием омега-квадрат} называется $\displaystyle S:=\left\{\omega ^{2}  >u_{1-\alpha }\right\}$, где $\displaystyle u_{1-\alpha }$ -- $\displaystyle ( 1-\alpha )$-квантиль распределения из предыдущего утверждения.
\end{definition}
\begin{definition}
    Критерии Пирсона, Колмогорова и омега-квадрат называют \textit{критериями согласия}, так как они проверяют гипотезу $\displaystyle H_{0} :\ P=P_{0}$ против альтернативы $\displaystyle H_{1} :\ P\neq P_{0}$.
\end{definition}